% !TeX spellcheck = en_US
\documentclass[french]{yLectureNote}

\title{Électromagnétisme 3}
\subtitle{Physique}
\author{Paulhenry Saux}
\date{\today}
\yLanguage{Français}
%jesse.groenen@cemes.fr
\professor{M.Brut}
\usepackage{graphicx}%----pour mettre des images
\usepackage[utf8]{inputenc}%---encodage
\usepackage{geometry}%---pour modifier les tailles et mettre a4paper
%\usepackage{awesomebox}%---pour les boites d'exercices, de pbq et de croquis ---d\'esactiv\'e pour les TP de PC
\usepackage{tikz}%---pour deiffner + d\'ependance de chemfig
% \usepackage{tabularx}%---pour dimensionner automatiquement les tableaux avec variable X
\usepackage{awesomebox}%---Pour les boites info, danger et autres
\usepackage{menukeys}%---Pour deiffner les touches de Calculatrice
\usepackage{fancyhdr}%---pour les en-t\^ete personnalis\'ees
\usepackage{blindtext}%---pour les liens
\usepackage{hyperref}%---pour les liens (\`a mettre en dernier)
\usepackage{caption}%---pour la francisation de la l\'egende table vers Tableau
\usepackage{pifont}
\usepackage{array}%---pour les tableaux
\usepackage{lipsum}
\usepackage{yFlatTable}
\usepackage{multicol}
\usepackage{tabularx}
\usepackage{booktabs}
\newcommand{\Lim}[1]{\lim\limits_{\substack{#1}}\:}
\renewcommand{\vec}{\overrightarrow}
\newcommand{\N}[0]{\mathbb{N}}
\newcommand{\dd}{\mathrm{d}}
\newcommand{\norm}[1]{||\vec{#1}||}
\newcommand{\fo}{\psi(\vec{r},t)}
\newcommand{\foe}{\psi(\vec{r},t)\*}
\newcommand{\HH}{\hat{H}}
\newcommand{\hb}{\hbar}
\newcommand{\lap}{\nabla^2}
\newcommand{\lapcc}{\frac{\partial^2 }{\partial x^2}+\frac{\partial^2 }{\partial y^2}+\frac{\partial^2 }{\partial z^2}}
\newcommand{\E}{\vec{E}(M)}
\newcommand{\B}{\vec{B}(M)}
\newcommand{\V}{V(M)}
\newcommand{\er}{\vec{e_{\rho}}}
\newcommand{\ep}{\vec{e_{\varphi}}}
\newcommand{\pip}{\pi^+}
\newcommand{\pim}{\pi^-}
\newcommand{\mv}{\mathcal{V}}
\DeclareMathOperator{\Div}{Div}
\newcommand{\rot}{\vec{\mathrm{rot}}}
\newcommand{\grad}{\vec{\mathrm{grad}}}
%\setcounter{chapter}{3}
\begin{document}
%\chapter{Puissance et énergie du champ Em}
%\setcounter{chapter}{4}

\chapter{Introduction}
\section{Échelles et propriétés}
Distances entre voisins :
\begin{itemize}
    \item phase dense : $L_a = 10^{-10}$
    \item phase gazeuse : Supérieur : $d>L_a$
\end{itemize}
Les milieux solides sont souvent électriquement neutre, à échelle plus ou moins grande selon s'il sont ordonnés. Elle est assurée à l'échelle $L_a$.

\begin{definition}[Polarisation]
    Le système est soumis à une excitation électrique
\end{definition}
\begin{definition}[Aimentation]
    Le système est soumis à une excitation magnétique
\end{definition}
\section{Moments dipolaires}
On étudie un ensemble des charges, l'ensemble est de dimension $a$. On utilise le principe de superposition.

\begin{flalign*}
    V(M) &= \sum \frac{q_i}{4\pi\epsilon_0}\frac{1}{||\vec{R}-\vec{r_i}||}\\
    &= \sum \frac{1}{R}(1+\frac{\vec{r_i}\cdot\vec{e_r}}{R}+\dots)\\
    &= \frac{Q}{4\pi\epsilon_0}\frac{1}{R}+ \frac{1}{4\pi\epsilon_0}\frac{1}{R^2}\vec{p}\cdot \vec{e_r}+\dots\\
    &= V_0(\vec{R})+V_1(\vec{R})+V_2(\vec{R})+\dots
\end{flalign*}
On a utilisé un DL pour obtenir la deuxième expression. %TODO Exliciter les étapes intermédiaires

On défini le moment dipolaire total : 
$$\vec{p} = \sum_i q_i\vec{ri}$$
\warningInfo{Approximation dipolaire}{On ne garde que $V_0$ et $V_1$}
En utilisant le gradient et avec ces approximations\marginInfo{Ce résultat est valide en statique mais aussi dans l'ARQS. En effet, on a $\delta t = \frac{R}{c}<<\tau$}, on obtient :
\begin{theorem}
    $$\vec{E} = \frac{Q}{4\pi\epsilon_0}\frac{1}{R^2}\vec{e_r}+\frac{Q}{4\pi\epsilon_0}\frac{3(\vec{p}\cdot \vec{e_r})\cdot \vec{e_r}-\vec{p}}{R^3}$$
\end{theorem}
\warningInfo{Édifice neutre}{Si Q=0, le premier terme, le terme dipolaire et $\vec{p}$ ne dépend pas de l'origine $O$.}
\includegraphics[scale=0.5]{s1}
\begin{definition}[Moment dipolaire magnétique]
    $$\vec{\mu} = \frac12\sum_i q_i (\vec{r_i}\wedge \vec{v_i}) = I\vec{S}$$ pour un courant circulaire de surface S
\end{definition}
\begin{theorem}[Champ magnétique loin des sources]
    $$\vec{B} = \frac{\mu_0}{4\pi}\frac{3(\vec{mu}\cdot \vec{e_r})\vec{e_r}-\vec{\mu}}{\vec{R}}$$
\end{theorem}
\includegraphics[scale=0.3]{s2}
\section{Passage du microscopique au macroscopique}
On n'a pas besoin de décrire les éléments à l'échelle microscopique, ce qui est très complexe. Cependant, les instruments de mesure ne permettent d'accéder qu'à une valeur moyenne locale, relative à un élément de volume $\delta v$ petit devant le volume $V$ du système mais grande devant $L_a^3$.

On va donc niveler les variations à l'échelle $\delta v$.

Le champ sera donc une valeur moyenne sur un volume petit\marginInfo{Comme les éuations de maxwell sont linéaires, elles s'appliquent très bien à ces valeur moyennées.}.

Le champ macroscopique est défini comme le champ moyennée sur une région suffisamment grande pour contenir assez de constituants élémentaites pur lisser les fluctutations microscipiques sans interets car non mesurables et suffisamment réduire pour assurer de ne pas perdre les variations mesurables à l'échelle de l'objet de volume V.

On a donc :
$$L_a^3\lll \delta v\lll  V$$ Il faut garder en tete que les valeurs de $\rho$ sont de valeurs moyennes.

On utilise donc 2 valeurs macroscopiques : $\vec{P}$ pour les milieux polarisés et $\vec{M}$ pour les milieux aimantés.
\end{document}
